% Capitulo: Dados para treinar IA
\chapter{Dados para treinar IA}
\label{cap:treinar-ia}

\rascunho{
  \textbf{Ideia do capítulo.} Quando o consumidor do dado é um modelo que vai aprender
  com ele, a modelagem passa a ser sobre exemplos: quais entram, como são rotulados,
  como são versionados e documentados, e o que nunca pode entrar. A maioria das
  empresas não vai treinar um modelo do zero, mas quase todas vão ajustar modelos e,
  principalmente, avaliar se a IA responde certo, e isso também é dado de treino.
  \begin{itemize}
    \item Diferença em relação ao capítulo~\ref{cap:modelagem-cd}: lá o dado é tabular
      e o modelo é nosso; aqui entram texto, imagem, pares de pergunta e resposta, e
      muitas vezes o modelo é de terceiros.
    \item Ligação com segurança e privacidade (ver \veja{sec:seguranca}).
  \end{itemize}
}

% ---------------------------------------------------------------------------
\section{Quando o consumidor é um modelo}
\label{sec:consumidor-modelo}

\rascunho{
  \begin{itemize}
    \item Treinar do zero, ajustar (\emph{fine-tuning}) e avaliar: três necessidades de
      dado diferentes, em ordem crescente de frequência nas empresas.
    \item O conjunto de avaliação (perguntas com a resposta certa conhecida) é o
      conjunto de treino que todo time de dados vai precisar montar, mesmo sem treinar
      nada.
  \end{itemize}
}

\pendente{escrever}

% ---------------------------------------------------------------------------
\section{Curadoria, rotulagem e versionamento}
\label{sec:curadoria}

\rascunho{
  \begin{itemize}
    \item Poucos exemplos bons valem mais que muitos ruins: duplicados, contraditórios
      ou fora do domínio estragam o resultado.
    \item Rotulagem: quem rotula, com que critério, e como medir a concordância entre
      rotuladores. \pendente{achar referência sobre rotulagem}
    \item Versionar o conjunto como se versiona código: saber exatamente com qual versão
      do dado cada modelo foi treinado (ver o capítulo~\ref{cap:dados-codigo}).
    \item Separar treino, validação e teste antes de qualquer coisa, e nunca deixar o
      teste vazar para o treino.
  \end{itemize}
}

\pendente{escrever}

% ---------------------------------------------------------------------------
\section{Documentar o conjunto de dados}
\label{sec:datasheets}

\rascunho{
  \begin{itemize}
    \item A proposta das fichas técnicas de conjuntos de dados (\emph{datasheets for
      datasets}): por que o conjunto foi criado, como foi coletado, o que contém, para
      que não deve ser usado \parencite{gebru2021}.
    \item A ficha é o equivalente, para dados, do que a documentação do modelo
      dimensional era para o BI.
  \end{itemize}
}

\pendente{escrever}

% ---------------------------------------------------------------------------
\section{Privacidade, direitos e viés}
\label{sec:privacidade-vies}

\rascunho{
  \begin{itemize}
    \item Dado pessoal no treino é difícil de tirar depois: anonimizar ou pseudonimizar
      antes, lembrando que dado pseudonimizado continua sendo dado pessoal pela LGPD.
      \pendente{conferir o texto da LGPD sobre dado pseudonimizado e citar a lei}
    \item Direitos sobre o conteúdo usado para treinar e termos de uso das fontes.
    \item Representatividade: um conjunto que só tem exemplos de uma região, um produto
      ou um perfil de cliente ensina o modelo a errar nos outros.
  \end{itemize}
}

\pendente{escrever}

% ---------------------------------------------------------------------------
\begin{leituras}
  \leitura{gebru2021}{A proposta de documentar cada conjunto de dados como se documenta
    um componente eletrônico: origem, composição, usos recomendados e limites.}
  \leitura{huyen2022}{Os capítulos sobre dados de treino tratam amostragem, rotulagem e
    desbalanceamento de forma prática.}
\end{leituras}
